% =============================================================
%  Design alternatives — topologies and devices that shorten the
%  recovery: full-bridge sub-modules, DC circuit breaker, fixed no-load
%  voltage, fast disconnect switches. Principles and trade-offs; the
%  decision taken on one project is in pugalur-thrissur.tex.
% =============================================================
\subsection{Design Alternatives}
\label{sec:dcf-alternatives}

The recovery times of the half-bridge symmetrical monopole are set by
mechanical operations: the AC circuit breaker, the bypass and discharge
switches, the tap-changer and the DC disconnect switches. A substantial
reduction requires a converter topology or a device that makes some of these
operations unnecessary~[24]. Four alternatives are known; each trades
investment or operating losses against recovery time.

\begin{itemize}[itemsep=0.2em]
  \item \textbf{Full-bridge sub-modules.} A full-bridge MMC is an open
        circuit when blocked, so the AC circuit breaker need not open and
        the AC network is less disturbed, and it operates at a modulation
        index above one, so it recovers to a reduced direct voltage without
        the tap-changer. Both remove mechanical operations from the sequence
        and shorten the recoveries and changeovers. The cost is a higher
        investment and higher operating losses, and fewer bidders offer it.
  \item \textbf{DC circuit breaker.} Clears the fault in the DC circuit and
        may keep the converters in service. Precautions for the DC
        overvoltage and the cable inrush remain and enter the achievable
        recovery time. The cost is the investment and, again, the number of
        bidders.
  \item \textbf{Fixed no-load voltage.} A converter-side no-load voltage
        chosen for the minimum direct voltage at the maximum modulation
        index allows recovery to any direct voltage within the time of the
        first restart, without tap-changer movement. The converter then
        operates at a higher AC current at rated voltage, with operating
        losses about 20\,\% higher~[24].
  \item \textbf{Fast disconnect switches.} Reduce the changeover to STATCOM
        operation to the few seconds of the fault clearing, but only when
        the link was already at reduced voltage, since otherwise the
        tap-changer sets the time. Their price is 5 to 10 times that of
        slow disconnects~[24].
\end{itemize}

\FloatBarrier
